% QMB_n_symbole.tex — Anhang: Formelzeichen und Abkürzungen
\chapter*{Formelzeichen und Abkürzungen}
\addcontentsline{toc}{chapter}{Formelzeichen und Abkürzungen}

\section*{Grundgrößen der FFGFT}

\begin{description}[leftmargin=4.5em,style=sameline,itemsep=0.4em]
  \item[$\xi$] Geometrieparameter der FFGFT, $\xi = 4/30000 =
    1/7500 \approx 1{,}333\times10^{-4}$. Einziger freier Parameter;
    folgt aus Galois-Gruppenordnungen des Trägers.
  \item[$\tilde T$] Kompakte Zeitkoordinate einer Mode; $\tilde T =
    1/m$ in natürlichen Einheiten ($\hbar = c = 1$).
  \item[$T^4/\mathbb{Z}_3$] FFGFT-Träger: vierdimensionaler Torus mit
    $\mathbb{Z}_3$-Symmetrie. Drei räumliche und eine zeitliche
    Wicklungsrichtung; neun Fixpunkte, Trialitätssymmetrie.
  \item[$K_{\mathrm{frak}}$] Fraktaler Korrekturfaktor,
    $K_{\mathrm{frak}} = 1 - 100\xi \approx 0{,}9867$. Entsteht beim
    Typ-I-Übergang vom kompakten Träger zur SI-Skala; kumulativ über
    17 Größenordnungen (RG-Lauf von Planck- zu elektroschwacher Skala).
  \item[$D_f$] Fraktale Raumdimension, $D_f = 3 - \xi \approx
    2{,}9999$. Lokal; nicht verwechseln mit $K_{\mathrm{frak}}$.
  \item[$L_0$] Fundamentale FFGFT-Länge, $L_0 = \xi\,\ell_P \approx
    2{,}15\times10^{-39}$\,m.
  \item[$\ell_P$] Planck-Länge, $\ell_P \approx 1{,}616\times10^{-35}$\,m.
  \item[$t_P$] Planck-Zeit, $t_P \approx 5{,}39\times10^{-44}$\,s.
\end{description}

\section*{Qubit- und Quantenrechner-Symbole}

\begin{description}[leftmargin=4.5em,style=sameline,itemsep=0.4em]
  \item[$(z,r,\theta)$] Zylindrische Phasenraumkoordinaten eines
    FFGFT-Qubits: $z\in[-1,+1]$ Projektion auf Berechnungsachse,
    $r\in[0,1]$ radiale Superpositionsamplitude, $\theta\in[0,2\pi)$
    relative Phase.
  \item[$(\alpha,\beta)$] Standard-QM-Amplituden eines Qubits:
    $|\psi\rangle = \alpha|0\rangle+\beta|1\rangle$,
    $|\alpha|^2+|\beta|^2 = 1$.
  \item[$K_{\mathrm{frak}}$ (Gatter)] In den $\pi$-Gattern
    ($\sigma_x$-Drehung): Drehwinkel $\alpha = \pi K_{\mathrm{frak}}$
    statt $\pi$. Geometrisch aus $D_f$ abgeleitet; prüfbare Vorhersage.
  \item[$\hat F$] Fundamentaler fraktaler Spektraloperator auf
    $\mathcal{H} = L^2(T^4/\mathbb{Z}_3)\otimes\mathbb{C}^3$.
    Selbstadjungiert; $\xi = \lambda_{\min}(\hat F_{D_4})$.
  \item[$D_4$] Wurzelgitter in $\mathbb{R}^4$: alle ganzzahligen
    Vektoren mit gerader Koordinatensumme. Kusszahl 24; Grundlage des
    Modenspektrums.
  \item[$\mathrm{CHSH}$] Clauser-Horne-Shimony-Holt-Parameter;
    lokal-realistisch $\le2$, Tsirelson-Schranke (QM)
    $2\sqrt2\approx2{,}828$.
  \item[$\Delta\mathrm{CHSH}_{\mathrm{elem}}$] Elementare
    FFGFT-Abweichung pro Messung: $\xi/(2\pi)\approx2\times10^{-5}$.
    Hypothese; formale Herleitung ausstehend.
  \item[$T_1,\,T_2$] Relaxationszeit ($T_1$) und Kohärenzzeit ($T_2$)
    eines Qubits; stets $T_2\le2T_1$.
  \item[$Q,\,Q_\varphi$] Register-Dimension: $Q = 2^n$ (Standard-QFT),
    $Q_\varphi = \varphi^n$ ($\phi$-QFT). Beide äquivalent für
    Periodenfindung bis $N < 2^{20}$ auf $\pm\xi$.
\end{description}

\section*{Spektraltheorie und Gitter}

\begin{description}[leftmargin=4.5em,style=sameline,itemsep=0.4em]
  \item[$N(n)$] Entartung der Schale $|k|^2 = n$ im $D_4$-Gitter;
    Koeffizient der Thetareihe
    $\theta_{D_4} = (\theta_3^4+\theta_4^4)/2$.
  \item[$Z_{\mathbb{Z}^4}(s)$] Spektral-Zeta des kubischen Torus:
    $8(1-4^{1-s})\zeta(s)\zeta(s-1)$.
  \item[$\mathrm{Aut}(D_4)$] Automorphismengruppe des $D_4$-Gitters;
    $|\mathrm{Aut}(D_4)| = 1152 = 2^7\cdot3^2$. Teilerfremd zu 5:
    keine Fünffachsymmetrie.
  \item[$\chi_0,\chi_1,\chi_2$] Die drei $\mathbb{Z}_3$-Charaktere des
    Orbifolds; jede $D_4$-Schale zerfällt gleichmäßig in
    $N(n)/3$ Zustände pro Sektor.
  \item[$\lambda_4$] Bogenlänge auf dem Massenkreis; de-Broglie-Relation
    $\lambda_4\cdot m = 2\pi$ exakt pro Mode (Ortsrelation, nicht
    Zeitrelation).
\end{description}

\section*{Zeitstruktur}

\begin{description}[leftmargin=4.5em,style=sameline,itemsep=0.4em]
  \item[Typ I] Messakt (Ausrollen): verlustbehaftet, $S^1_m\to\mathbb{R}_t$.
    Informationsverlust, Zeitpfeil geometrisch.
  \item[Typ III] Darstellungswechsel (Pullback): verlustfrei, bijektiv.
    Kein Informationsverlust; $K_{\mathrm{frak}}$ nicht beteiligt.
  \item[$\xi_{\mathrm{Higgs}}$] Higgs-Kopplungsparameter,
    $\approx 1{,}038\times10^{-5}$. Wirkt im algebraischen
    Zustandsraum (Kopplungsraten, Dekohärenz); \emph{nicht} identisch
    mit $\xi$, der im geometrischen Ortsraum wirkt.
\end{description}

\section*{Statusmarker}

\begin{description}[leftmargin=2.5em,style=sameline,itemsep=0.3em]
  \item[\textbf{[K]}] Konfirmiert — vollständig aus $\xi$ oder einer
    korpusgesicherten Ableitung bewiesen; Prüfskript bestanden.
  \item[\textbf{[B]}] Belegt — numerisch gesichert; analytischer Beweis
    vorhanden oder trivial.
  \item[\textbf{[S]}] Spekulativ — Hypothese oder Hinweis; noch nicht
    vollständig hergeleitet.
  \item[\textbf{[E]}] Etablierte Physik — Standardergebnis, hier nur
    in Erinnerung gerufen.
  \item[\textbf{[X]}] Geprüft negativ — explizit ausgeschlossen mit
    Begründung.
  \item[\textbf{[SETZUNG]}] Axiomatische Setzung — Ausgangspunkt der
    Ableitung, nicht weiter begründet.
\end{description}
